\subsection{Transformers}

The Transformer architecture \cite{vaswani_attention_2017} revolutionized sequence modeling by introducing an attention-based approach that eliminates the need for recurrence or convolution operations.

The transformer architecture consists of three main components. First, the input processing stage converts raw input (such as text) into a format suitable for the model often through tokenization and an embedding projection. Second, multiple transformer blocks process this information through self-attention and feed-forward mechanisms. Finally, the output processing stage maps the transformed representations to the desired output format through a linear layer followed by a softmax activation, producing probability distributions over possible outputs.

\subsubsection{Input processing}
For text processing, the input undergoes tokenization - a process that segments text into units called tokens. These tokens can represent individual characters, subwords, or whole words, depending on the tokenization strategy. Each token is mapped to a unique integer index, which is then transformed into a high-dimensional vector through an embedding layer. The embedding process transforms each token index $i$ into a dense vector through a learned embedding matrix $E$:
\begin{equation}
\text{embedding}(i) = E_i
\end{equation}

where $E \in \mathbb{R}^{|V|\times d} $ is the embedding matrix, $|V|$ is the vocab size, and $ d $ is the embedding dimension.

Additionally, positional encodings are added to these embeddings to preserve sequence order information, since the self-attention mechanism itself is permutation-invariant. The positional encoding PE is added to these embeddings to inject position information. For position $po$s and dimension $d_{\text{model}}$, the encoding is computed as:
\begin{align}
PE_{(pos,2i)} &= \sin\left(\frac{pos}{10000^{2i/d_{\text{model}}}}\right) \\
PE_{(pos,2i+1)} &= \cos\left(\frac{pos}{10000^{2i/d_{\text{model}}}}\right)
\end{align}

where $pos$ is the position in the sequence and $i$ is the embedding index. The final input representation $X$ for each position $pos$ is the sum of its embedding and positional encoding:

\begin{equation}
X_{pos} = \text{embedding}(token_{pos}) + PE_{pos}
\end{equation}

\subsubsection{Self-attention}
At its core, the architecture computes self-attention through scaled dot-product attention:

\begin{equation}
\text{Attention}(Q, K, V) = \text{softmax}\left(\frac{QK^\top}{\sqrt{d_k}}\right)V \label{eq:attention}
\end{equation}

where $Q$, $K$, and $V$ represent learned query, key, and value matrices respectively, and $d_k$ is the dimension of the key vectors. The scaling factor $\sqrt{d_k}$ prevents the dot products from growing too large in magnitude, which could push the softmax function into regions with extremely small gradients. This mechanism enables each position to attend to all other positions, facilitating the modeling of long-range dependencies while supporting parallel computation.

The attention mechanism can be divided into projection, calculation of scores, softmax, and output. First the input sequence $X$ is projected to a query $Q$, key $K$, and value $V$ matrix by multiplying with the correspond weight matrix:

\begin{equation}
    Q = XW_Q, \quad K = XW_K, \quad V = XW_V
\end{equation}


The attention scores are then computed through matrix multiplication of Q and K, followed by scaling and then softmax:
\begin{equation}
    \text{Scores}_{ij} = \frac{\exp\left(\frac{(QK^\top)_{ij}}{\sqrt{d_k}}\right)}{\sum_{k=1}^n \exp\left(\frac{(QK^\top)_{ik}}{\sqrt{d_k}}\right)}
\end{equation}

These scores represent the relative importance of all other positions for each position in the sequence. The final attention output is obtained by multiplying these scores with $V$ to get a weighted sum of values:

\begin{equation}
\text{Output} = \text{Scores} \times V
\end{equation}

In practice multi-head attention is used, allowing the model to jointly attend to information from different representation subspaces:

\begin{equation}
\text{MultiHead}(Q, K, V) = \text{Concat}(\text{head}_1, \ldots, \text{head}_h)W^O
\end{equation}

where each head computes attention (\ref{eq:attention}) with different learned projections of queries, keys, and values. This multi-head mechanism enables the model to capture various types of dependencies and patterns in the data simultaneously. The outputs of all heads are concatenated and linearly transformed by the output weight matrix $W^O$ to produce the final attention output for that layer.

Each transformer block consists of two main sub-layers: multi-head attention followed by a position-wise feed-forward network, with each sub-layer wrapped in a residual connection and layer normalization. The computations can be expressed as:
\begin{align}
& \text{AttentionBlock}(X) = \text{LayerNorm}(X + \text{MultiHead}(X, X, X)) \\
& \text{TransformerBlock}(X) = \text{LayerNorm}(\text{AttentionBlock}(X) + \text{FFN}(\text{AttentionBlock}(X)))
\end{align}

The position-wise feed-forward network (FFN) applies two linear transformations with a ReLU activation in between:

\begin{equation}
\text{FFN}(x) = \text{max}(0, xW_1 + b_1)W_2 + b_2
\end{equation}

where $W_1 \in \mathbb{R}^{d_{\text{model}} \times d_{\text{ff}}}$, $W_2 \in \mathbb{R}^{d_{\text{ff}} \times d_{\text{model}}}$, and $b_1, b_2$ are learnable parameters. The dimensionality $d_{\text{ff}}$ is typically larger than $d_{\text{model}}$, allowing the network to learn richer feature transformations.

Layer normalization is applied to each sub-layer's output and computes the mean $\mu(x)$ and variance $\sigma(x)^2$ across the feature dimension:

\begin{equation}
\text{LayerNorm}(x) = \gamma \left(\frac{x - \mu}{\sqrt{\sigma^2 + \epsilon}}\right) + \beta
\end{equation}

where $H$ is the hidden dimension, $\gamma$ and $\beta$ are learnable parameters, and $\epsilon$ is a small constant for numerical stability.
The complete flow through a transformer block can thus be written as:

\begin{align}
Z_1 &= X + \text{MultiHead}(X, X, X) \\
Z_2 &= \text{LayerNorm}(Z_1) \\
Z_3 &= Z_2 + \text{FFN}(Z_2) \\
\text{output} &= \text{LayerNorm}(Z_3)
\end{align}